% TODO make hw stylesheet
\documentclass[10pt]{article}
\usepackage[letterpaper,top=1in,left=1in,right=1in]{geometry}
\usepackage[dvipsnames,svgnames]{xcolor}
\usepackage[shortlabels]{enumitem}
\usepackage[framemethod=TikZ]{mdframed}
\usepackage{amsmath,amssymb,amsthm}
\usepackage[colorlinks]{hyperref}
\usepackage{multicol}
\usepackage{tabularx}
\usepackage{bbm}

\hypersetup{
    colorlinks=true,
    linkcolor=blue,
    filecolor=magenta,
    urlcolor=magenta,
    pdftitle={MAT 535 HW}
}

\usepackage{mathtools}
\usepackage{thmtools}
\usepackage[textsize=footnotesize]{todonotes}
\usepackage{titling}
\usepackage{cancel}

\reversemarginpar
\mdfdefinestyle{mdbluebox}{roundcorner=10pt,innerbottommargin=9pt,
linecolor=blue,backgroundcolor=TealBlue!5,}
\declaretheoremstyle[headfont=\sffamily\bfseries\color{MidnightBlue},
mdframed={style=mdbluebox},]{thmbluebox}

\mdfdefinestyle{mdbloobox}{frametitlefont=\bfseries,innerbottommargin=8pt,
nobreak=true,backgroundcolor=TealBlue!5,linecolor=blue,}
\declaretheoremstyle[headfont=\bfseries\color{MidnightBlue},
mdframed={style=mdbloobox},headpunct={\\[3pt]},postheadspace=0pt,]{thmbloobox}

\mdfdefinestyle{mdredbox}{frametitlefont=\bfseries,innerbottommargin=8pt,
nobreak=true,backgroundcolor=Salmon!5,linecolor=RawSienna,}
\declaretheoremstyle[headfont=\bfseries\color{RawSienna},
mdframed={style=mdredbox},headpunct={\\[3pt]},postheadspace=0pt,]{thmredbox}

\mdfdefinestyle{mdgreenbox}{linecolor=ForestGreen,backgroundcolor=ForestGreen!5,
linewidth=2pt,rightline=false,leftline=true,topline=false,bottomline=false,}
\declaretheoremstyle[headfont=\bfseries\sffamily\color{ForestGreen!70!black},
mdframed={style=mdgreenbox},headpunct={ --- },]{thmgreenbox}

\mdfdefinestyle{mdorangebox}{linecolor=RedOrange,backgroundcolor=RedOrange!5,
linewidth=2pt,rightline=false,leftline=true,topline=false,bottomline=false,}
\declaretheoremstyle[headfont=\bfseries\sffamily\color{RedOrange!70!black},
mdframed={style=mdorangebox},headpunct={ --- },]{thmorangebox}

\mdfdefinestyle{mdblackbox}{linecolor=black,backgroundcolor=RedViolet!5!gray!5,
linewidth=3pt,rightline=false,leftline=true,topline=false,bottomline=false,}
\declaretheoremstyle[mdframed={style=mdblackbox}]{thmblackbox}

\declaretheoremstyle[spaceabove=6pt,spacebelow=6pt]{boldhead}
\declaretheoremstyle[spaceabove=6pt,spacebelow=6pt,headfont=\normalfont\itshape]{italichead}

\declaretheorem[style=boldhead,name=Problem]{problem}
\declaretheorem[style=boldhead,name=Problem,numbered=no]{problem*}
\declaretheorem[style=boldhead,name=Setup]{setup} % for config geo
\declaretheorem[style=boldhead,name=Example,sibling=problem]{example}
\declaretheorem[style=boldhead,name=$\bigstar$ Problem,sibling=problem]{niceprob}
\declaretheorem[style=boldhead,name=Theorem,sibling=problem]{theorem}
\declaretheorem[style=boldhead,name=Corollary,sibling=theorem]{corollary}
\declaretheorem[style=boldhead,name=Proposition,sibling=theorem]{prop}
\declaretheorem[style=boldhead,name=Lemma,numberwithin=problem]{lemma}
\declaretheorem[style=boldhead,name=Theorem,numbered=no]{theorem*}
\declaretheorem[style=boldhead,name=Lemma,numbered=no]{lemma*}
\declaretheorem[style=boldhead,name=Claim,numberwithin=problem]{claim}
\declaretheorem[style=boldhead,name=Claim,numbered=no]{claim*}
\declaretheorem[style=boldhead,name=Remark,numberwithin=problem]{remark}
\declaretheorem[style=boldhead,name=Remark,numbered=no]{remark*}
\declaretheorem[style=boldhead,name=Definition,sibling=theorem]{definition}
\declaretheorem[style=boldhead,name=Definition,numbered=no]{definition*}

\providecommand{\ol}{\overline}
\providecommand{\ceil}[1]{\left\lceil #1 \right\rceil}
\providecommand{\floor}[1]{\left\lfloor #1 \right\rfloor}
\newcommand{\dd}{\,\mathrm{d}}
\providecommand{\ul}{\underline}
\providecommand{\eps}{\varepsilon}
\providecommand{\half}{\frac{1}{2}}
\providecommand{\CC}{\mathbb C}
\providecommand{\FF}{\mathbb F}
\providecommand{\II}{\mathbb I}
\providecommand{\NN}{\mathbb N}
\providecommand{\QQ}{\mathbb Q}
\providecommand{\RR}{\mathbb R}
\providecommand{\ZZ}{\mathbb Z}
\providecommand{\ind}{\mathbbm 1}
\providecommand{\dg}{^\circ}
\providecommand{\ii}{\item}
\providecommand{\setdiff}{\smallsetminus}
\providecommand{\alert}[1]{{\sffamily\textbf{\textcolor{blue}{#1}}}}
\providecommand{\inv}{^{-1}}
\providecommand{\mb}[1]{\mathbf{#1}}
\newcommand{\what}{\widehat}

\newenvironment{soln}{\begin{proof}[Solution]}{\end{proof}}

\providecommand{\pow}{\operatorname{Pow}}
\providecommand{\vocab}[1]{{\textbf{\textcolor{ForestGreen}{#1}}}}
\providecommand{\lcm}{\operatorname{lcm}}
\providecommand{\abs}[1]{\left\lvert #1\right\rvert}
\providecommand{\smabs}[1]{\lvert #1\rvert}
\providecommand{\innerprod}[1]{\langle\!\langle #1\rangle\!\rangle}
\providecommand{\norm}[1]{\left\| #1\right\|}
\providecommand{\smnorm}[1]{\| #1\|}
\providecommand{\gr}{\operatorname{gr}}

\usepackage{mathrsfs}

% place title at top right
\setlength{\droptitle}{-8ex}
\pretitle{\begin{flushright}\Large\bfseries}
\posttitle{\par\end{flushright}}
\preauthor{\begin{flushright}}
\postauthor{\end{flushright}}
\predate{\begin{flushright}}
\postdate{\end{flushright}}

\title{MAT 535 HW01}
\author{\large Aditya Pahuja}
\date{\large Due 02/05}
\begin{document}
\maketitle

\newcommand{\Hx}{\mathcal H{\times}}
\newcommand{\Kx}{\mathcal K{\times}}
\begin{problem}
    [Appendix II \S1 \#2]
    Let $H$ be a group. Define a map $\Hx$
    from $\mathbf{Grp}$ to itself on objects and morphisms as follows:
    \begin{gather*}
        \Hx(G) = H\times G,\qquad
	\Hx(\varphi\colon G_1\to G_2) =
	\{(h,g)\mapsto(h,\varphi(g))\}.
    \end{gather*}
    Prove that $\Hx$ is a functor.
\end{problem}

\begin{soln}
    We only need to check that for any group $G$,
    $\Hx(\operatorname{id}_G) = \operatorname{id}_{H\times G}$
    and, for homomorphisms $\varphi\colon G_1\to G_2$
    and $\psi\colon G_2\to G_3$, $\Hx(\psi\circ\varphi)
    = \Hx(\psi)\circ\Hx(\varphi)$. The former is true because
    $\Hx(\operatorname{id}_G)$ is the map sending
    $(h,g)\in H\times G$ to $(h,\operatorname{id}_G(g)) = (h,g)$,
    and the latter is true because for any $(h,g)\in H\times G_1$,
    \begin{equation*}
        \Hx(\psi)\circ\Hx(\varphi)(h,g)
	= \Hx(\psi)(h,\varphi(g))
	= (h,\psi\circ\varphi(g))
	= \Hx(\psi\circ\varphi)(h,g).\qedhere
    \end{equation*}
\end{soln}

\begin{problem}
    [Appendix II \S1 \#3]
    Show that the map $\mathbf{Ring}$ to $\mathbf{Grp}$
    by mapping a ring to its group of units defines a functor.
    Show by explicit examples that this functor is neither faithful nor full.
\end{problem}

\begin{soln}
    Consider the map from $\mathbf{Ring}$ to $\mathbf{Grp}$
    sending each ring $R$ to its group of units $R^\times$ and
    sending the ring homomorphism $\varphi\colon R_1\to R_2$
    to the group homomorphism
    $\varphi|_{R_1^\times}\colon R_1^\times\to R_2^\times$,
    which is well-defined because if $r\in R_1$ is a unit then
    $\varphi(r)\varphi(r\inv) = 1_{R_2}$ implies that
    $\varphi(r)$ is also a unit. This is a functor because
    the identity map on a ring $R$ of course restricts to
    the identity on $R^\times$, and given two ring homomorphisms
    $\varphi\colon R_1\to R_2$ and $\psi\colon R_2\to R_3$,
    \begin{equation*}
	(\psi\circ\varphi)|_{R_1^\times}
	= \psi\circ(\varphi|_{R_1^\times})
	= \psi|_{R_2^\times}\circ\varphi|_{R_1^\times}
    \end{equation*}
    where we are using the fact that
    $\varphi|_{R_1^\times}\subseteq R_2^\times$ to get the second equality.

    Denote the functor above by $\mathcal F$.
    To show that this functor is neither full nor faithful,
    consider the map
    \begin{equation*}
	\operatorname{Hom}_{\mathbf{Ring}}(\ZZ[x],\ZZ)\overset{\mathcal F}\to
	\operatorname{Hom}_{\mathbf{Grp}}(\{1,-1\},\{1,-1\}).
    \end{equation*}
    This map sends every homomorphism $\varphi\colon\ZZ[x]\to\ZZ$
    to the identity on $\{1,-1\}$ because $\varphi(1) = 1$
    implies $\varphi(-1) = -1$; since there is more than one homomorphism
    $\varphi\colon \ZZ[x]\to\ZZ$ (there is a bijection between integers $n$
    and homomorphisms $\varphi$ with $\varphi(x) = n$),
    we find that $\mathcal F$ is not faithful. Of course, the map sending
    both $1$ and $-1$ to $1$ is also a group homomorphism,
    so $\mathcal F$ is not full, either.
\end{soln}

\pagebreak

\begin{problem}
    [Appendix II \S2 \#2]
    Let $H$ and $K$ be groups and let $\Hx$ and $\Kx$
    be the functors on $\mathbf{Grp}$ defined in the previous section.
    Let $\varphi\colon H\to K$ be a group homomorphism.
    \begin{enumerate}[(a)]
        \ii Show that the maps $\eta_A\colon H\times A\to K\times A$
	given by $\eta_A(h,a) = (\varphi(h),a)$ determine
	a natural transformation $\eta$ from $\Hx$ to $\Kx$.
	\ii Show that the transformation $\eta$ is a natural isomorphism
	if and only if $\varphi$ is a group isomorphism.
    \end{enumerate}
\end{problem}

\begin{soln}
    (a) To show that $\eta$ is a natural transformation, we need to show that
    \begin{equation*}
        \eta_B\circ\Hx(\psi) = \Kx(\psi)\circ\eta_A
    \end{equation*}
    for all group homomorphisms $\psi\colon A\to B$.
    This follows from
    \begin{equation*}
	\eta_B(\{\Hx(\psi)\}(h,a))
	= \eta_B(h,\psi(a)) = (\varphi(h),\psi(a))
	= \{\Kx(\psi)\}(\varphi(h),a) = \{\Kx(\psi)\}(\eta_A(h,a))
    \end{equation*}
    for all $(h,a)\in H\times A = \Hx(A)$.
    \\

    (b) If $\eta$ is a natural isomorphism for all groups $A$, then,
    taking $A$ to be the trivial group,
    $\eta_1\colon H\times 1\to K\times 1$
    is an isomorphism. Since the projections
    $p_H\colon H\times 1\to H$ and $p_K\colon K\times 1\to K$
    onto the first coordinate are isomorphisms, we get that
    \begin{equation*}
        \varphi = p_K\circ\eta_1\circ p_H\inv\colon H\to K
    \end{equation*}
    is an isomorphism, where the equality comes from
    \begin{equation*}
        p_K\circ\eta_1\circ p_H\inv(h) =
	p_K(\eta_1(h,1)) = p_K(\varphi(h),1) = \varphi(h).
    \end{equation*}

    Conversely, if $\varphi$ is an isomorphism, then
    $\eta_A(h,a) = (\varphi(h),a)$ has an inverse
    $\eta_A\inv(k,a) = (\varphi\inv(k),a)$, so each $\eta_A$ is an isomorphism,
    meaning $\eta$ is a natural isomorphism of $\Hx$ and $\Kx$.
\end{soln}

\begin{problem}
    [Appendix II \S2 \#3]
    Express the universal property of the commutator quotient group
    as a universal arrow for some functor $\mathcal F$.
\end{problem}

\begin{soln}
    Let $\mathcal F$ be the functor from $\mathbf{Ab}$ to $\mathbf{Grp}$
    sending each abelian group and homomorphism of abelian groups
    to itself. For a group $G$, let $U(G) = G/G'$
    and $\iota\colon G\to \mathcal FU(G) = U(G)$ be the projection homomorphism.
    Then, for any abelian group $H$ and group homomorphism
    $\varphi\colon G\to \mathcal FH=H$, the universal property of
    the commutator quotient group dictates that there is a unique
    morphism of abelian groups $\Phi\colon G/G'\to H$
    such that $\mathcal F(\Phi)\iota = \varphi$, so
    $(U(G),\iota)$ is a universal arrow from $X$ to
    the ``inclusion functor'' $\mathcal F$.
\end{soln}










\end{document}
